% Generated draft: requires mathematical and source review.
\documentclass[11pt]{article}
\usepackage[margin=25mm]{geometry}
\usepackage{fontspec}
\setmainfont{Latin Modern Roman}
\usepackage{amsmath,amssymb,mathtools,unicode-math}
\setmathfont{Latin Modern Math}
\usepackage{enumitem,needspace,xurl,hyperref}
\setlength{\parindent}{0pt}
\setlength{\parskip}{5pt}
\setlength{\emergencystretch}{3em}
\newcommand{\N}{\mathbb{N}}
\newcommand{\Z}{\mathbb{Z}}
\newcommand{\Q}{\mathbb{Q}}
\newcommand{\R}{\mathbb{R}}
\newcommand{\C}{\mathbb{C}}
\newcommand{\F}{\mathbb{F}}
\newcommand{\A}{\mathbb{A}}
\newcommand{\PP}{\mathbb{P}}
\newcommand{\E}{\mathbb{E}}
\newcommand{\eps}{\varepsilon}
\newcommand{\dd}{\,\mathrm{d}}
\begin{document}
\section*{UnsolvedMath ID: 701662}
\textbf{Source problem code:} catalogue entry
\subsection{Definitions and mathematical statement}
Let \(G\) be a finite group. For \(x\in G\), let \(\langle x\rangle\) denote the cyclic subgroup generated by \(x\). The group \(G\) is called \emph{cut}, or \emph{inverse semi-rational}, if
\[
\forall x,y\in G,\qquad
\langle x\rangle=\langle y\rangle
\ \Longrightarrow\ 
\exists g\in G\;\bigl(g^{-1}xg=y\ \text{or}\ g^{-1}xg=y^{-1}\bigr).
\]
Let \(\operatorname{Irr}(G)\) be the set of ordinary irreducible complex characters of \(G\), and define the field generated by its ordinary character-table entries by
\[
\Q(G)=\Q\bigl(\chi(g):\chi\in\operatorname{Irr}(G),\ g\in G\bigr)\subseteq\C.
\]
For a prime \(p\), write \(\operatorname{O}_p(G)\) for the largest normal subgroup of \(G\) whose order is a power of \(p\). For a finite group \(H\), its exponent \(\exp(H)\) is the least positive integer \(n\) such that \(h^n=1\) for every \(h\in H\); in particular, \(\exp(\{1\})=1\). A group is solvable if its derived series reaches the trivial subgroup.

\begin{enumerate}
\item[(a)] Does there exist a real constant \(c>0\), independent of \(G\), such that
\[
[\Q(G):\Q]\leq c
\qquad\text{for every finite cut group }G?
\]
The supplied source reports that such a uniform bound exists when \(G\) is additionally assumed solvable, attributing this result to J. F. Tent, \emph{J. Algebra} \textbf{363} (2012), 73--82. This is retained as a source-reported qualification, without independent verification.

\item[(b)] Is every Sylow \(3\)-subgroup of every finite cut group itself cut? Explicitly, if \(P\leq G\) has order equal to the largest power of \(3\) dividing \(|G|\), must
\[
\forall x,y\in P,\qquad
\langle x\rangle=\langle y\rangle
\ \Longrightarrow\ 
\exists u\in P\;\bigl(u^{-1}xu=y\ \text{or}\ u^{-1}xu=y^{-1}\bigr)
\]
hold?

\item[(c)] For every finite solvable cut group \(G\) and every \(p\in\{5,7\}\), is it true that
\[
\exp\bigl(\operatorname{O}_p(G)\bigr)\mid p?
\]
\end{enumerate}

\emph{Source qualifications.} The local record labels this multipart problem open; that label is supplied registry information, not an independently verified assessment of any part. The record also refers to additional information and positive results without specifying which parts they concern. Its proposer attribution is corrupted, and one journal volume designation is unclear; these source defects are retained explicitly in the references and require review.
\subsection{Short English statement}
For finite cut groups, ask whether the field generated by all ordinary character-table entries has uniformly bounded degree over \(\Q\), and whether every Sylow \(3\)-subgroup is cut. For finite solvable cut groups, also ask whether \(\operatorname{O}_p(G)\) has exponent dividing \(p\) for \(p=5,7\). The supplied source reports a positive answer to the first question under the additional solvability assumption. Its open label and literature claims have not been independently verified.
\subsection{Sources}
\begin{itemize}
\item Local UnsolvedMath registry, permanent ID \(701662\), exact source problem code \texttt{catalogue entry}, titled \emph{Kourovka Notebook Problem 20.4}. The supplied record belongs to \emph{Kourovka Notebook (Issues 1--21)} and carries the status \texttt{open}. Its creation and update timestamps are both \texttt{2026-10-16T14:00:00Z}; these are registry timestamps, not dates of independent literature verification. The supplied upstream revision is \texttt{d371ee917ac9647fbb06cd6fc529deec6f3dda82}.
\item \emph{The Kourovka Notebook (Unsolved Problems in Group Theory)}, original problem 20.4. The supplied background identifies the 20th issue (2022) and the 21st issue edition (2026), with identifier \texttt{arXiv:1401.0300v48}. The set description names E. I. Khukhro and V. D. Mazurov as editors. These bibliographic details are supplied source information and have not been independently checked.
\item The supplied proposed year is 2022. The proposer field is truncated as \texttt{A.\ B\{\textbackslash{}"a}; the background repeats this fragment and contains a separated fragment \texttt{chle\}} after its discussion heading. The complete proposer name cannot be reliably transcribed from this corrupted attribution and has not been reconstructed.
\item J. F. Tent, \emph{J. Algebra} \textbf{363} (2012), 73--82. The supplied statement cites this reference for the uniform degree bound in part (a) under the additional assumption that \(G\) is solvable. No article title is supplied, and the result has not been independently verified here.
\item The supplied statement refers to additional information and some positive results in \emph{Adv. Group Theory Appl.}, volume designation \texttt{8B} as supplied, 2020, 157--160, or \url{https://arxiv.org/abs/2001.02637}. No author or article title is supplied for this reference, and the volume designation and scope of the positive results require checking.
\end{itemize}
\end{document}
